\section{把人、钱和信息接进帝国}
\label{sec:han-imperial-capacity}

长安能够发出诏令，不等于诏令已经抵达每一户人家。西汉的制度史最容易被官名和法令填满，真正困难的是追问这些名目怎样经过郡县、王国、商人、宗族与边塞吏员，才成为可用的粮、兵和消息。前一三四年举孝廉与前一一九年前后的财政整顿，是两种不同的接口：前者把地方评价送入选官，后者将钱币、盐铁和运输纳入中枢可见的账目。它们都增加了朝廷的视野，却没有消除谁替国家奔走、谁从中获利的问题。

从现存材料看，西汉地方并不是一片同质的腹地。关中接近皇帝和仓储；旧楚地、王国、边郡、新附的南方与河西各自拥有不同的道路、首领和生产条件。征南越、入滇、灭卫满朝鲜之后的新郡，首先是任命、文书和军事上的安排，而不是当地社会在某一天突然变成关中式县乡。帝国以不同速度把不同地区接入，正是它能够长久扩展、又始终难以整齐控制的原因。

\subsection{尺度不同的材料}

前四九年以后，外戚王氏逐渐占据中枢位置。外戚并不是“制度失灵”后才出现的异常人物；幼主、婚姻、尚书、禁军和官爵分配使亲属网络本来就能成为权力的运输线。元始二年的户口数字尤其诱人，因为它似乎给出一个可以精确丈量的帝国。实际上，这是一份来自行政申报的尺度：它说明朝廷仍试图用郡国簿籍表达所掌握的人口与土地，却不能直接等于真实总人口，更不能推出每户税役的负担。

同样，钱币、城址和简牍让制度不再只是帝纪中的抽象名词。它们能校正某处铸造、一个关塞、一道道路或一个官署的物质条件；它们不能由局部材料填成全国性的统计图。承认这个限度并非缩小汉帝国。相反，它使“统一”成为一种反复取得、需要文书和中介不断维持的能力，而不是地图上一次完成的涂色。

\begin{causalchain}[从财政集中到政治风险]
远征与新郡经营增加军费和转运压力；财政与选官措施扩大中枢可见的资源；可被看见、可被调度的位置也更值得外戚、近臣、地方官和宗室争夺。这里不是一条单线因果，而是国家能力与竞争资源彼此放大的反馈回路。
\end{causalchain}

这种视角也改变了我们读“治世”的方式。轻税、稳定币制或吏治整顿可以是真实的政策选择，却不是一张能覆盖所有地区、所有阶层的社会成绩单。地方文书留下得很不均衡，正史又偏向记录朝廷如何理解自己。西汉最可贵的成就，或许是能在如此大的差异中不断重建接口；它最危险的脆弱性，也恰在接口被少数集团垄断、或被远方压力挤断之时。
